\documentclass[aspectratio=169]{beamer}
\usepackage{iftex}
\ifXeTeX
  \usepackage{fontspec}
  \setsansfont{DejaVu Sans}
\else
  \usepackage[utf8]{vietnam}
\fi
\usepackage{amsmath}
\usepackage{amssymb}
\usepackage{graphicx}
\usepackage{booktabs}

% Cài đặt theme
\usetheme{Madrid}
\usecolortheme{default}

% Thông tin trang bìa
\title[Looped Flows \& Flow Map LMs]{Looped Flows \& Flow Map Language Models}
\subtitle{Cải thiện Suy luận Lặp và Mô hình hóa Ngôn ngữ Đơn bước}
\author{Hung Manh}
\date{\today}

\begin{document}

% Slide Title
\begin{frame}
    \titlepage
\end{frame}

% Slide Outline
\begin{frame}{Nội dung trình bày}
    \tableofcontents
\end{frame}

%---------------------------------------------------------
\section{Tóm tắt}
\begin{frame}{Tóm tắt}
    \begin{itemize}
        \item \textbf{Vấn đề:} Giải bài toán khó cần nhiều compute time. Các looped models thực hiện điều này qua việc cập nhật tuần hoàn hidden state. Tuy nhiên, huấn luyện gradient ngắn hạn khiến việc học các early updates gặp khó khăn.
        \item \textbf{Đề xuất:} \textbf{Looped flows} - huấn luyện recurrency với các mục tiêu khử nhiễu cục bộ.
        \item \textbf{Cơ chế:} 
        \begin{itemize}
            \item \textit{Temporal alignment}: Giảm dần mức độ nhiễu và sử dụng shared noise.
            \item \textit{Inference}: Tích phân vận tốc của probability flow kết hợp với các looped states.
        \end{itemize}
        \item \textbf{Kết quả:} Vượt qua các SOTA looped models, đạt \textbf{58.8\%} trên ARC-AGI-1 và \textbf{12.2\%} trên ARC-AGI-2; cho phép scale tính toán và tạo multiple valid predictions.
    \end{itemize}
\end{frame}

%---------------------------------------------------------
\section{Giới thiệu}
\begin{frame}{1. Giới thiệu: Bối cảnh \& Thách thức}
    \textbf{Bối cảnh \& Động lực:}
    \begin{itemize}
        \item Giải bài toán khó đòi hỏi gia tăng tính toán lúc inference.
        \item Khác với Autoregressive LLMs, các looped models sử dụng distributed state và cập nhật tuần hoàn để tích lũy computation.
    \end{itemize}
    \vspace{0.3cm}
    \textbf{Thách thức của Looped Models:}
    \begin{itemize}
        \item Huấn luyện bằng full BPTT tốn kém bộ nhớ và dễ gặp hiện tượng vanishing/exploding gradients.
        \item Thực tế chỉ backprop qua vài bước $\rightarrow$ early updates không nhận được tín hiệu để hỗ trợ future computation.
        \item Hậu quả: Chuỗi lặp không ổn định, sinh ra các failure modes ở cả bài toán đơn giản.
    \end{itemize}
\end{frame}

\begin{frame}{1. Giới thiệu: Giải pháp Looped Flows}
    \textbf{Đóng góp chính của Looped Flows:}
    \begin{itemize}
        \item Học một chuỗi mục tiêu khử nhiễu cục bộ theo thời gian ở nhiều mức nhiễu khác nhau.
        \item \textbf{Temporal alignment:} Phân bổ huấn luyện theo thời gian, giảm dần mức nhiễu, sử dụng shared noise $\rightarrow$ hidden states duy trì tính hữu ích xuyên suốt.
        \item \textbf{Cơ chế Inference:} Giải bài toán bằng cách tích phân trường vận tốc của probability flow kết hợp với recurrent states.
        \item \textbf{Lợi ích:}
        \begin{itemize}
            \item Scale compute qua một finer time grid.
            \item Đưa ra nhiều dự đoán hợp lệ bằng stochastic/Euler integration.
            \item Học được stable recurrency, khắc phục điểm yếu của looped models cũ.
        \end{itemize}
    \end{itemize}
\end{frame}

%---------------------------------------------------------
\section{Kiến thức nền tảng}
\begin{frame}{2. Kiến thức nền: Looped Models}
    \begin{itemize}
        \item \textbf{Định nghĩa:} Cập nhật hidden state $z \in \mathbb{R}^{d_z}$ qua hàm $f$:
        $$z_i = f(z_{i-1}; c)$$
        \item Đưa ra dự đoán qua đầu đọc $g$: $\hat{x}_i = g(z_i)$.
        \item \textbf{Hàm loss thực tế:} Để tránh full BPTT, mô hình dùng toán tử dừng gradient $sg(\cdot)$:
        $$L_{loop}(f, g) := \mathbb{E}_{c,x} \left[\sum_i CE(\hat{x}_i, x)\right] \quad \text{với } z_i = f(sg(z_{i-1}); c)$$
        \item \textit{Điểm yếu:} Tổn thất tại bước $\ge i$ không truyền tín hiệu về quá khứ $z_{<i}$. Quá trình lặp deterministic chỉ cho phép 1 lời giải duy nhất.
    \end{itemize}
\end{frame}

\begin{frame}{2. Kiến thức nền: Flow and Diffusion Models}
    \begin{itemize}
        \item \textbf{Continuous flow matching:} Đi từ nhiễu $x_0 \sim \mathcal{N}(0, \sigma^2 I)$ đến dữ liệu thật $x_1 \sim p_{data}$.
        \item \textbf{Probability path:} Nội suy $I_t := (1-t)x_0 + t x_1$.
        \item \textbf{Velocity field:} $\dot{x}_t = b_t(x_t; c)$.
        \item Tích phân số (VD: Euler tiến) để xấp xỉ lời giải:
        $$x_{t_{i+1}} = x_{t_i} + (t_{i+1} - t_i) \hat{b}_{t_i}(x_{t_i}; c)$$
        \item \textbf{Denoising function ($D_t$):} Thay vì dự đoán trực tiếp vận tốc, học hàm khử nhiễu độc lập tại mỗi $t$:
        $$L_{flow}(\hat{D}) := \mathbb{E}[CE(\hat{x}_t, x_1)]$$
        \item \textit{Nhược điểm:} Khó học các tính toán tuần tự do phụ thuộc hoàn toàn vào flow.
    \end{itemize}
\end{frame}

%---------------------------------------------------------
\section{Phương pháp luận}
\begin{frame}{3. Looped Flows: Cơ chế Huấn luyện}
    \textbf{Stateful Denoiser:}
    \begin{itemize}
        \item Khử nhiễu $x_{t_i}$ đồng thời cập nhật looped state $z_{t_i}$:
        $$(\hat{x}_{t_i}, z_{t_i+1}) = \hat{D}_{t_i}(I_{t_i}, sg(z_{t_i}); c)$$
    \end{itemize}
    
    \textbf{Temporal Alignment:}
    \begin{itemize}
        \item Sắp xếp lưới thời gian $t_0 < t_1 < \dots < t_k$ để mức nhiễu giảm dần.
        \item \textbf{Shared noise:} Mẫu $x_0$ dùng chung cho toàn bộ quỹ đạo $\rightarrow$ ép mô hình liên kết các thời điểm, giúp $z_{t_i}$ mang thông tin hữu ích cho tương lai.
    \end{itemize}
    
    \textbf{Adaptive Computation Time (ACT):}
    \begin{itemize}
        \item Dùng đầu đọc phân loại nhị phân $q(z_{t_i+1})$ để quyết định dừng suy luận sớm khi độ chính xác đã bão hòa.
    \end{itemize}
\end{frame}

\begin{frame}{3. Looped Flows: Cơ chế Suy luận}
    \begin{itemize}
        \item \textbf{Stochastic Integration:} 
        \begin{itemize}
            \item Thêm hệ số ngẫu nhiên $\gamma \ge 0$.
            \item Lùi nhẹ $x_{t_i}$ về thời điểm $s < t_i$ (tiêm nhiễu mới) trước khi thực hiện bước Euler tiến. 
            \item Khi $\gamma = 0$, mô hình quay về Euler ODE truyền thống.
        \end{itemize}
        \item \textbf{Inference-time Scaling:}
        \begin{itemize}
            \item Giải bài toán bằng finer time grid so với lúc huấn luyện.
            \item Sử dụng kỹ thuật \textbf{best-Q ensembling}: Chạy nhiều mẫu, chọn dự đoán có xác suất thành công cao nhất từ đầu đọc ACT $q$.
        \end{itemize}
    \end{itemize}
\end{frame}

%---------------------------------------------------------
\section{Nghiên cứu liên quan}
\begin{frame}{4. Nghiên cứu liên quan}
    \begin{itemize}
        \item \textbf{Looped models (TRM, FPRM):} Thiếu giám sát toàn cục do cắt gradient, dễ rơi vào spurious attractors. \textit{Looped flows giải quyết bằng cách khởi tạo lặp từ local denoising.}
        \item \textbf{Stochastic looped models:} Tận dụng nhiễu nhưng vẫn phụ thuộc gradient cục bộ. \textit{Looped flows kết hợp lợi ích ngẫu nhiên và giải quyết triệt để học cục bộ.}
        \item \textbf{Flow \& Diffusion Models:} Học khử nhiễu độc lập. \textit{Looped flows dùng local objectives tạo lộ trình học mượt mà cho hidden state.}
        \item \textbf{Energy models:} \textit{Looped flows tương đương việc học chuỗi năng lượng làm mịn giảm dần}, giúp quá trình tối ưu dễ dàng hơn.
    \end{itemize}
\end{frame}

%---------------------------------------------------------
\section{Thử nghiệm và Đánh giá}
\begin{frame}{5. Thử nghiệm: Chính xác \& Khả năng mở rộng}
    \textbf{So sánh đơn quỹ đạo (Single-trajectory):}
    \begin{itemize}
        \item Looped flows áp đảo TRM, FPRM, GRAM trên:
        \begin{itemize}
            \item Sudoku-Extreme (\textbf{97.9\%})
            \item Maze-Hard (\textbf{86.7\%})
            \item ARC-AGI-1 (\textbf{58.8\%} - tăng từ 44.6\% của TRM)
            \item ARC-AGI-2 (\textbf{12.2\%} - tăng từ 7.8\% của TRM)
        \end{itemize}
    \end{itemize}
    \textbf{Inference-time scaling:}
    \begin{itemize}
        \item Ensembling cải thiện thêm hiệu năng (Sudoku 99.3\%, ARC 59.5\%).
        \item Finer time grid tăng accuracy mạnh: Sudoku từ 74.5\% (8 steps) $\rightarrow$ 97.9\% (128 steps).
    \end{itemize}
    \textbf{Độ ổn định:}
    \begin{itemize}
        \item Giải quyết 90.9\% các trường hợp thất bại của TRM (bao gồm non-convergence và spurious attractors).
    \end{itemize}
\end{frame}

\begin{frame}{5. Thử nghiệm: Đa dạng Lời giải \& Ablation Study}
    \textbf{Đa dạng lời giải (Diversity):} Đánh giá trên bài toán đa nghiệm.
    \begin{itemize}
        \item N-Queens (8x8): \textbf{99.9\%} ACC / \textbf{91.4\%} Coverage.
        \item Graph Coloring (8 đỉnh): \textbf{0.7} conflicts / \textbf{89.4\%} Coverage.
        \item Vượt trội hoàn toàn về khả năng tìm ra tập lời giải đa dạng nhờ probability transport.
    \end{itemize}
    
    \textbf{Ablation Study:}
    \begin{itemize}
        \item \textbf{Flow:} Bỏ time conditioning, bỏ interpolant, không giảm nhiễu, hay không chia sẻ nhiễu đều làm giảm mạnh hiệu năng.
        \item \textbf{Recurrence:} So với mô hình không lặp (FLM) hoặc self-conditioning, Looped flows đạt khả năng tổng quát hóa tốt nhất trên test set.
        \item Tích phân SDE ($\gamma=5$) cho kết quả tốt hơn Euler ODE ($\gamma=0$).
    \end{itemize}
\end{frame}

%---------------------------------------------------------
\section{Kết luận}
\begin{frame}{6. Kết luận}
    \begin{itemize}
        \item \textbf{Looped Flows} là một framework thanh lịch, giải quyết vấn đề huấn luyện của Looped Models thông qua các mục tiêu khử nhiễu cục bộ.
        \item \textbf{Temporal alignment} giúp trạng thái lặp xây dựng computation hữu ích trên toàn cục, khắc phục hoàn toàn điểm yếu của gradient ngắn hạn.
        \item Mở ra khả năng \textbf{inference-time scaling} mạnh mẽ qua lưới thời gian mịn và tích phân ngẫu nhiên.
        \item \textbf{Hướng tương lai:} Phát triển các thuật toán huấn luyện không cần mô phỏng (simulation-free training).
    \end{itemize}
\end{frame}

% thêm bài báo khác vào chỗ này
% ==============================================================================
% File: elf_paper.tex
% Bài báo: Embedded Language Flows (ELF)
% Mô hình khuếch tán ngôn ngữ liên tục dựa trên Flow Matching
% ==============================================================================

\section{Embedded Language Flows (ELF)}

% Slide 0: Giới thiệu tổng quan bài báo
\begin{frame}{ELF: Tổng quan bài báo}
    \begin{itemize}
        \item \textbf{Tên bài báo:} \textit{Embedded Language Flows (ELF)}
        \item \textbf{Ý tưởng cốt lõi:} Mô hình khuếch tán ngôn ngữ liên tục dựa trên Flow Matching trong không gian nhúng biểu diễn (latent space).
        \item \textbf{Đóng góp chính:}
        \begin{itemize}
            \item Chứng minh continuous diffusion hoàn toàn vượt trội discrete diffusion trên ngôn ngữ nếu khử nhiễu hoàn toàn trong không gian nhúng liên tục.
            \item Chỉ thực hiện rời rạc hóa (discretization/decoding) tại bước cuối cùng ($t = 1$).
            \item Sử dụng mạng chia sẻ trọng số giữa denoising và decoding, tích hợp sẵn Classifier-Free Guidance (CFG) lúc huấn luyện.
        \end{itemize}
    \end{itemize}
\end{frame}

% Slide 1: Bối cảnh & Phân loại
\begin{frame}{1. ELF: Bối cảnh \& Phân loại DLMs}
    \textbf{Bối cảnh:}
    \begin{itemize}
        \item Diffusion và flow-based models là tiêu chuẩn vàng trong sinh ảnh/video, đang được mở rộng sang xử lý ngôn ngữ.
    \end{itemize}
    \vspace{0.2cm}
    \textbf{Phân loại mô hình khuếch tán ngôn ngữ (DLMs):}
    \begin{itemize}
        \item \textit{Rời rạc (Discrete):} Thao tác trực tiếp trên token (trước đây được ưu tiên do dễ tối ưu hóa).
        \item \textit{Liên tục (Continuous):} Ánh xạ token sang không gian biểu diễn liên tục để tiến hành khử nhiễu.
    \end{itemize}
    \vspace{0.2cm}
    \textbf{Giải pháp đề xuất của ELF:}
    \begin{itemize}
        \item Khử nhiễu hoàn toàn trong không gian nhúng liên tục.
        \item Định nghĩa trường vận tốc qua thời gian liên tục với Flow Matching.
        \item Dùng mạng chia sẻ trọng số, chỉ rời rạc hóa tại bước $t = 1$.
    \end{itemize}
\end{frame}

% Slide 2: Khung làm việc ELF
\begin{frame}{2. ELF: Khung làm việc (Framework)}
    \begin{itemize}
        \item \textbf{Từ rời rạc sang liên tục:} Chuỗi token $s \in \mathcal{V}^L$ được chuyển sang biểu diễn liên tục $x$ bằng bộ mã hóa T5 đông cứng.
        \item \textbf{Nội suy tuyến tính (Rectified Flows):} Tạo biến ẩn $z_t$ dọc theo đường đi:
        $$z_t = t x + (1 - t)\epsilon, \quad t \in [0, 1], \quad \epsilon \sim \mathcal{N}(0, \mathbf{I})$$
        \item \textbf{Trường vận tốc dòng (Flow Velocity):} Đạo hàm theo thời gian:
        $$v = \frac{dz}{dt} = x - \epsilon$$
        \item \textbf{Dự đoán $x$ (x-prediction):} Trực tiếp dự đoán biểu diễn sạch $x_\theta(z_t, t)$ thay vì dự đoán vận tốc $v$, giúp mạng nơ-ron hoạt động ổn định ở số chiều cực cao.
    \end{itemize}
\end{frame}

% Slide 3: Hàm mục tiêu & Rời rạc hóa
\begin{frame}{3. ELF: Hàm mục tiêu \& Rời rạc hóa}
    \begin{itemize}
        \item \textbf{Nhánh khử nhiễu (Denoising - chiếm 80\% trọng số):} Tối thiểu hóa sai số MSE của trường vận tốc dự đoán:
        $$\mathcal{L}_{\text{MSE}} = \mathbb{E}_{t, x, \epsilon} \left[ \frac{1}{(1-t)^2} \|x_\theta(z_t, t) - x\|^2 \right]$$
        \item \textbf{Nhánh giải mã (Decoding - chiếm 20\% trọng số):}
        \begin{itemize}
            \item Tại bước cuối $t = 1$, biến đổi $z_t$ thành phân phối token qua ma trận unembedding $W$.
            \item Dùng chung trọng số với mạng khử nhiễu, huấn luyện bằng Cross-Entropy:
            $$\mathcal{L}_{\text{CE}} = \mathbb{E}_{\tilde{z}} \left[ \text{CrossEntropy}(W x_\theta(\tilde{z}), s) \right]$$
        \end{itemize}
    \end{itemize}
\end{frame}

% Slide 4: Điều kiện hóa & Dẫn hướng
\begin{frame}{4. ELF: Điều kiện hóa \& Dẫn hướng (Guidance)}
    \begin{itemize}
        \item \textbf{Self-conditioning:} Thực hiện lượt truyền thứ hai điều kiện hóa trên dự đoán trung gian $\hat{x}'$ từ bước trước để tăng độ mượt của quỹ đạo.
        \item \textbf{Classifier-Free Guidance (CFG):} Dễ dàng áp dụng trong miền liên tục:
        $$v_{\text{cfg}}(z_t \mid c) = \omega v(z_t \mid c) + (1 - \omega) v(z_t \mid \emptyset)$$
        \item \textbf{Training-time CFG:} Tích hợp thẳng CFG vào hàm mục tiêu vận tốc ngay từ lúc huấn luyện $\rightarrow$ loại bỏ hoàn toàn chi phí tính toán kép khi suy luận.
        \item \textbf{Sinh có điều kiện:} Ghép prompt vào trước chuỗi mục tiêu; phần prompt được giữ nguyên không làm nhiễu.
    \end{itemize}
\end{frame}

% Slide 5: Kết quả thực nghiệm
\begin{frame}{5. ELF: Kết quả thực nghiệm}
    \textbf{Sinh không điều kiện (OpenWebText):}
    \begin{itemize}
        \item Ở 32 bước lấy mẫu, ELF-B (105M) đạt \textbf{Gen. PPL 24.08}, vượt trội so với các đối thủ rời rạc như MDLM (104.29) và Duo (62.47).
        \item \textbf{Hiệu quả dữ liệu cực cao:} Đạt đỉnh chỉ với 45B token huấn luyện (ít hơn $\sim$10 lần so với $>500$B tokens của các baseline).
    \end{itemize}
    \vspace{0.2cm}
    \textbf{Chưng cất lũy tiến (Progressive Distillation):}
    \begin{itemize}
        \item Cho phép giảm số bước lấy mẫu xuống mức cực ít (\textbf{1 đến 4 bước}) mà vẫn duy trì chất lượng sinh văn bản xuất sắc.
    \end{itemize}
    \vspace{0.2cm}
    \textbf{Sinh có điều kiện:}
    \begin{itemize}
        \item \textbf{Tóm tắt XSum:} ROUGE-1 đạt \textbf{36.0} (đánh bại các mô hình tự hồi quy cùng kích cỡ).
        \item \textbf{Dịch máy WMT14:} Đạt BLEU \textbf{26.4}.
    \end{itemize}
\end{frame}

% Slide 6: Kết luận
\begin{frame}{6. ELF: Kết luận \& Đóng góp}
    \begin{itemize}
        \item \textbf{Khẳng định:} Mô hình khuếch tán ngôn ngữ liên tục hoàn toàn có thể cạnh tranh và vượt trội hơn các phương pháp rời rạc.
        \item \textbf{Bí quyết tối ưu của ELF:}
        \begin{enumerate}
            \item Giữ toàn bộ quỹ đạo làm nhiễu trong không gian nhúng liên tục.
            \item Chỉ thực hiện rời rạc hóa ở bước thời gian cuối cùng ($t = 1$).
            \item Sử dụng mạng chia sẻ trọng số, loại bỏ bộ giải mã phụ trợ rườm rà.
        \end{enumerate}
        \item \textbf{Tác động:} Mở ra tiềm năng tạo văn bản chất lượng cao với chi phí huấn luyện thấp và tốc độ suy luận cực nhanh.
    \end{itemize}
\end{frame}

% ==============================================================================
% File: fmlm_paper.tex
% Bài báo: Flow Map Language Models: One-step Language Modeling via Continuous Denoising
% Tác giả: Chanhyuk Lee, Nicholas M. Boffi, Jinwoo Kim et al. (KAIST & CMU, 2026)
% ==============================================================================

%---------------------------------------------------------
\section{FMLM: Tóm tắt}
\begin{frame}{FMLM: Tóm tắt bài báo}
    \begin{itemize}
        \item \textbf{Vấn đề:} Discrete diffusion models hứa hẹn sinh song song thay thế autoregressive LLMs, nhưng chất lượng sụp đổ thảm hại ở few-step regime ($1 \sim 4$ bước) do lỗi factorization approximation.
        \item \textbf{Đề xuất:} \textbf{Flow Map Language Models (FMLM)} — mô hình hóa ngôn ngữ bằng continuous flows trên biểu diễn one-hot unconstrained.
        \item \textbf{Cơ chế:}
        \begin{itemize}
            \item \textit{Flow Language Model (FLM)}: Đưa bài toán về denoiser trên simplex, huấn luyện bằng cross-entropy và time reparameterization theo decoding error.
            \item \textit{Two-time denoiser ($\delta_{s,t}$)}: Tham số hóa flow map nằm trọn trên simplex, chưng cất bằng semigroup condition.
        \end{itemize}
        \item \textbf{Kết quả:} FMLM đạt khả năng sinh văn bản \textbf{1 bước (one-step)} vượt qua chất lượng 8 bước của discrete diffusion baselines, tăng tốc \textbf{8.3$\times$}, hỗ trợ continuous guidance ổn định.
    \end{itemize}
\end{frame}

%---------------------------------------------------------
\section{FMLM: Bối cảnh \& Thách thức}
\begin{frame}{FMLM: Giới hạn của Discrete Diffusion}
    \textbf{Nút thắt của Autoregressive LLMs:}
    \begin{itemize}
        \item Sinh tuần tự từng token một ($O(L)$ bước tính toán), tạo ra điểm nghẽn lớn về tốc độ khi sinh chuỗi dài.
    \end{itemize}
    \vspace{0.2cm}
    \textbf{Tử huyệt Factorization Error của Discrete Diffusion:}
    \begin{itemize}
        \item Discrete diffusion (MDLM, Duo) khử nhiễu trên không gian chuỗi $\mathcal{V}^L$.
        \item Do không gian trạng thái quá lớn, mô hình bắt buộc phải dùng xấp xỉ phân rã tích:
        $$\tilde{p}_{t|s}(\mathbf{y}_t | \mathbf{y}_s) := \prod_{l=1}^L p^l_{t|s}(y_t^l | \mathbf{y}_s)$$
        \item Giả định các token độc lập có điều kiện chỉ đúng trong giới hạn vô cùng bé ($t \to s$).
        \item \textit{Hậu quả ở few-step ($1 \sim 4$ bước):} Bước nhảy lớn làm mất hoàn toàn tương quan ngữ nghĩa giữa các token $\rightarrow$ văn bản vô nghĩa (PPL $> 1000$) hoặc lặp từ liên tục (entropy collapse $< 4.0$).
    \end{itemize}
\end{frame}

\begin{frame}{FMLM: Bước chuyển tư duy (Paradigm Shift)}
    \textbf{Bác bỏ định kiến về Discrete Modalities:}
    \begin{itemize}
        \item Quan niệm phổ biến cho rằng dữ liệu rời rạc bắt buộc phải dùng discrete noising process.
        \item \textbf{Mệnh đề 3.5 (Giới hạn toán học):} Trong không gian rời rạc hữu hạn $S$, không tồn tại bất kỳ deterministic map $f: S \to S$ nào ở cấp độ mẫu để biến đổi một phân phối tùy ý về phân phối đích.
        \item Discrete diffusion không thể có sample-level deterministic flow map để nhảy thẳng từ nhiễu về dữ liệu trong 1 bước.
    \end{itemize}
    \vspace{0.2cm}
    \textbf{Giải pháp Continuous Flows:}
    \begin{itemize}
        \item Nhúng token dạng one-hot lossless vào không gian Euclid không ràng buộc $\mathbb{R}^{L \times |\mathcal{V}|}$.
        \item Dòng liên tục trong $\mathbb{R}^{L \times |\mathcal{V}|}$ sở hữu một flow map tất định duy nhất ở cấp độ mẫu, mở đường cho việc sinh 1 bước chính xác.
    \end{itemize}
\end{frame}

%---------------------------------------------------------
\section{FMLM: Mô hình cơ sở FLM}
\begin{frame}{FMLM: Mô hình cơ sở Flow LM (FLM)}
    \begin{itemize}
        \item \textbf{Stochastic Interpolant:} Nội suy đường đi giữa Gaussian noise $\mathbf{x}_0 \sim \mathcal{N}(0, \mathbf{I})$ và one-hot data $\mathbf{x}_1$:
        $$I_t := (1 - t)\mathbf{x}_0 + t\mathbf{x}_1$$
        \item \textbf{Thách thức dự đoán velocity:} Học vận tốc $\mathbf{b}_t = \mathbb{E}[\mathbf{x}_1 - \mathbf{x}_0 | I_t]$ bằng MSE bị rank bottleneck do $\mathbf{x}_0$ full-rank $\rightarrow$ underfitting nghiêm trọng.
        \item \textbf{Chuyển sang Denoiser trên Simplex (Bổ đề 3.1):}
        $$D_t(\mathbf{x}) := \mathbb{E}[\mathbf{x}_1 | I_t = \mathbf{x}], \quad \mathbf{b}_t(\mathbf{x}) = \frac{D_t(\mathbf{x}) - \mathbf{x}}{1 - t}$$
        \begin{itemize}
            \item $D_t(\mathbf{x})^l = p_{1|t}^l(\cdot | I_t = \mathbf{x}) \in \Delta^{|\mathcal{V}|-1}$ (luôn nằm trên simplex xác suất).
            \item Huấn luyện bằng \textbf{Cross-Entropy loss} thích ứng với hình học one-hot, chặn chặt khoảng cách Wasserstein $W_2^2 \le C\epsilon$ (Mệnh đề 3.2).
        \end{itemize}
    \end{itemize}
\end{frame}

\begin{frame}{FMLM: Tái tham số hóa thời gian (Time Reparameterization)}
    \begin{itemize}
        \item \textbf{Vấn đề:} Tỷ lệ lỗi giải mã $P_e(t) := \frac{1}{L}\sum_l \mathbb{P}(g^l(\mathbf{x}_t) \ne g^l(\mathbf{x}_1))$ dốc đứng sát $t = 1$.
        \begin{itemize}
            \item Lấy mẫu đều $t \sim \mathcal{U}[0, 1]$ lãng phí tín hiệu ở vùng chưa quyết định token và bỏ sót vùng quyết định then chốt khi $|\mathcal{V}|$ lớn.
        \end{itemize}
        \item \textbf{Đề xuất hàm $\tau(t)$ chuẩn hóa tiến độ giải mã:}
        $$\tau(t) = 1 - \frac{|\mathcal{V}|}{|\mathcal{V}| - 1} P_e(t)$$
        \item \textbf{Lợi ích:}
        \begin{itemize}
            \item Huấn luyện và lấy mẫu đều trong không gian $\tau(t) \sim \mathcal{U}[0, 1]$.
            \item Mỗi bước tính toán đóng góp đồng đều vào việc giảm decoding error.
            \item Giúp mô hình huấn luyện ổn định với từ vựng lớn $|\mathcal{V}| \approx 50.000$.
        \end{itemize}
    \end{itemize}
\end{frame}

%---------------------------------------------------------
\section{FMLM: Phương pháp luận Flow Map}
\begin{frame}{FMLM: Two-time Denoiser trên Simplex}
    \textbf{Mục tiêu Flow Map:}
    \begin{itemize}
        \item Ánh xạ trực tiếp giữa hai thời điểm bất kỳ: $X_{s,t}(\mathbf{x}) = \mathbf{x} + (t - s)\mathbf{v}_{s,t}(\mathbf{x})$.
        \item Cho phép sinh 1 bước qua $\hat{\mathbf{x}}_1 = X_{0,1}(\mathbf{x}_0)$.
        \item \textit{Thách thức:} $X_{s,t}$ và $\mathbf{v}_{s,t}$ rời khỏi simplex do nhiễu Gaussian $\rightarrow$ dùng MSE cho chất lượng rất kém.
    \end{itemize}
    \vspace{0.2cm}
    \textbf{Định nghĩa Two-time Denoiser ($\delta_{s,t}$):}
    $$\delta_{s,t}(\mathbf{x}) := \mathbf{x} + (1 - s)\mathbf{v}_{s,t}(\mathbf{x})$$
    \textbf{4 Tính chất cốt lõi (Mệnh đề 3.3):}
    \begin{enumerate}
        \item Phục hồi chính xác Flow Map: $X_{s,t}(\mathbf{x}) = \frac{1 - t}{1 - s}\mathbf{x} + \frac{t - s}{1 - s}\delta_{s,t}(\mathbf{x})$.
        \item Bảo toàn trên simplex: $\delta_{s,t}(\mathbf{x})^l \in \Delta^{|\mathcal{V}|-1}$ $\rightarrow$ dùng được softmax output.
        \item Khớp với denoiser thông thường trên đường chéo: $\delta_{s,s}(\mathbf{x}) = D_s(\mathbf{x})$.
        \item Thỏa mãn semigroup dạng tổ hợp lồi: $\delta_{s,t} = \gamma \delta_{s,u} + (1 - \gamma)\delta_{u,t}(X_{s,u})$.
    \end{enumerate}
\end{frame}

\begin{frame}{FMLM: Huấn luyện Chưng cất & Suy luận 1 bước}
    \textbf{Hàm mất mát Semigroup Distillation (Mệnh đề 3.4):}
    \begin{itemize}
        \item Tận dụng tính chất tổ hợp lồi trên simplex, tối ưu qua KL divergence:
        $$\mathcal{L}_{\text{KL}}(\hat{\delta}) = \mathbb{E}_{s,u,t} \left[ \sum_{l=1}^L \text{KL}(\bar{\delta}^l_{s,t} \parallel \hat{\delta}^l_{s,t}(I_s)) \right] + \mathbb{E}_{t} \left[ \sum_{l=1}^L \text{KL}(\hat{D}_t(I_t)^l \parallel \hat{\delta}^l_{t,t}(I_t)) \right]$$
        với teacher: $\bar{\delta}_{s,t} := sg\left(\gamma \hat{\delta}_{s,u}(I_s) + (1 - \gamma)\hat{\delta}_{u,t}(\hat{X}_{s,u})\right)$, và $\hat{D}_t$ là FLM teacher đã đóng băng.
    \end{itemize}
    \vspace{0.2cm}
    \textbf{Kỹ thuật then chốt:}
    \begin{itemize}
        \item \textbf{Boundary sampling:} Do $\tau(t)$ phẳng gần $t = 0$, cố định xác suất $p = 1/32$ lấy mẫu cặp $(s, t) = (0, 1)$ để bơm gradient mạnh cho bước nhảy 1-step.
        \item \textbf{One-step inference:} $\hat{\mathbf{x}}_1 = \hat{\delta}_{0,1}(\mathbf{x}_0)$ sau đó decode qua argmax.
    \end{itemize}
\end{frame}

%---------------------------------------------------------
\section{FMLM: Dẫn hướng trong suy luận}
\begin{frame}{FMLM: Continuous Inference Guidance}
    \textbf{Autoguidance (Nâng cao chất lượng sinh không điều kiện):}
    \begin{itemize}
        \item Ngoại suy vận tốc ra xa model yếu: $\hat{\mathbf{b}}_t^{\text{(guided)}} = \hat{\mathbf{b}}_t^{\text{(weak)}} + \eta (\hat{\mathbf{b}}_t - \hat{\mathbf{b}}_t^{\text{(weak)}})$.
        \item Nhờ hoạt động trong không gian Euclid, FLM duy trì tính ổn định ở $\eta$ lên đến \textbf{100} (giảm PPL từ 96.9 xuống \textbf{51.6} tại $\eta=50$).
        \item Discrete diffusion sụp đổ hoàn toàn ở $\eta \ge 10$ do ngoại suy logit nhân bản factorization error.
    \end{itemize}
    \vspace{0.2cm}
    \textbf{Reward Guidance (FMRG - Điều hướng theo thuộc tính):}
    \begin{itemize}
        \item Flow map $X_{t_n, 1}$ cung cấp khả năng \textbf{look-ahead vi phân} tới điểm cuối:
        $$\mathbf{x}_{t_{n+1}} = X_{t_n, t_{n+1}}(\mathbf{x}_{t_n}) + \lambda \nabla r\left( X_{t_n, 1}(\mathbf{x}_{t_n}) \right)$$
        \item Classifier/Reward model $r$ chỉ cần huấn luyện trên \textbf{văn bản sạch}, không cần train trên dữ liệu nhiễu như discrete diffusion.
    \end{itemize}
\end{frame}

%---------------------------------------------------------
\section{FMLM: Thử nghiệm và Đánh giá}
\begin{frame}{FMLM: Hiệu năng sinh nhiều bước \& Đột phá One-step}
    \textbf{Chế độ nhiều bước (1024 steps):}
    \begin{itemize}
        \item FLM vượt trội các baseline SOTA trên cả LM1B và OpenWebText (OWT):
        \begin{itemize}
            \item LM1B ($L=128$): FLM PPL \textbf{96.91} vs Duo 98.14, MDLM 109.21, CANDI 120.99.
            \item OWT ($L=1024$): FLM PPL \textbf{62.23} vs Duo 77.69, MDLM 105.15.
        \end{itemize}
    \end{itemize}
    \vspace{0.2cm}
    \textbf{Chế độ cực ít bước (1, 2, 4 steps) — FMLM One-step Breakthrough:}
    \begin{itemize}
        \item Discrete diffusion chưng cất (MDLM+SDTT, Duo+DCD, Di4C) sụp đổ ở 1 bước (PPL $> 1000$ hoặc entropy collapse $< 4.0$).
        \item \textbf{FMLM 1-step:} Đạt PPL \textbf{119.34} trên LM1B và \textbf{168.30} trên OWT.
        \begin{itemize}
            \item Chất lượng 1 bước tương đương mô hình rời rạc chưng cất ở $8 \sim 16$ bước.
            \item Tăng tốc độ sinh \textbf{8.3$\times$}.
            \item Self-BLEU (0.073) và GPT-4.1 Judge xác nhận không bị mode collapse.
        \end{itemize}
    \end{itemize}
\end{frame}

\begin{frame}{FMLM: Cấu trúc logic (Sudoku) \& Ablation Study}
    \textbf{Mô hình hóa ràng buộc logic toàn cục (Sudoku $9 \times 9$):}
    \begin{itemize}
        \item \textit{Unconditional 1-step:} FMLM đạt \textbf{9.38\%} lời giải hợp lệ (discrete baselines đều $= 0.00\%$).
        \item \textit{Conditional 1-step (20 clues):} FMLM giải đúng \textbf{39.84\%} chỉ trong 1 forward pass.
    \end{itemize}
    \vspace{0.2cm}
    \textbf{Ablation Study (Bảng 18):}
    \begin{itemize}
        \item \textbf{Loss \& Parameterization:} Denoiser + Softmax + CE (\textbf{96.91}) $\gg$ Denoiser + MSE (129.04) $\gg$ Velocity + MSE (3801.36 - không hội tụ).
        \item \textbf{Time Schedule:} Decoding error schedule $\tau(t)$ (\textbf{106.98}) $\gg$ Uniform (149.18).
        \item \textbf{Distillation Objective:} Semigroup (\textbf{119.34}) $\gg$ Eulerian (149.13) $\gg$ Lagrangian (193.08) do không cần JVP tốn kém và bất ổn định.
        \item \textbf{Boundary sampling ($p=1/32$):} Cải thiện PPL 1-step từ 142.61 xuống \textbf{119.34}.
    \end{itemize}
\end{frame}

%---------------------------------------------------------
\section{FMLM: Kết luận}
\begin{frame}{FMLM: Kết luận \& Định hướng}
    \begin{itemize}
        \item \textbf{Continuous Flow vượt trội trên ngôn ngữ:} Bác bỏ giả thuyết bắt buộc phải dùng discrete diffusion cho dữ liệu rời rạc; continuous flow cho chất lượng cao hơn và suy luận nhanh hơn.
        \item \textbf{Two-time Denoiser trên Simplex:} Cầu nối lý thuyết giúp chưng cất flow map ổn định bằng cross-entropy / semigroup condition.
        \item \textbf{Khả năng mở rộng:} Hỗ trợ autoguidance ổn định và reward-guided generation hiệu quả với look-ahead vi phân.
        \item \textbf{Hạn chế \& Hướng tương lai:}
        \begin{itemize}
            \item Ma trận nhúng $O(L|\mathcal{V}|)$ tốn thêm khoảng $30\%$ bộ nhớ lúc huấn luyện $\rightarrow$ hướng giải quyết bằng sparse gradients.
            \item Tiềm năng ứng dụng lớn cho reinforcement learning finetuning (RLFT) nhờ giảm chi phí rollout tính reward.
        \end{itemize}
    \end{itemize}
\end{frame}

% ==============================================================================
% File: lpt_paper.tex
% Bài báo: The Surprising Effectiveness of Shared Memory in Looped Transformers
% Tác giả: Giovanni Monea, Keshav Ramji, Yousef El-Kurdi, Luis A. Lastras,
%          Yoav Artzi, Nathan Godey, Ramón Fernandez Astudillo
%          (Cornell và IBM Research, arXiv:2610.02383, 2026)
% Nguồn số liệu: reports/2610.02383_shared_memory_looped_transformers.md
% ==============================================================================

\section{Bộ nhớ dùng chung trong Transformer vòng lặp LPT}

% ------------------------------------------------------------------------------
% Slide 0: Tổng quan bài báo
% ------------------------------------------------------------------------------
\begin{frame}{LPT: Tổng quan bài báo}
    \begin{itemize}
        \item \textbf{Tên bài báo:} \textit{The Surprising Effectiveness of Shared Memory in Looped Transformers}
        \item \textbf{Tác giả và tổ chức:} Giovanni Monea, Keshav Ramji, Yousef El-Kurdi et al., Cornell và IBM Research, arXiv:2610.02383
        \item \textbf{Ý tưởng cốt lõi:} Huấn luyện Transformer vòng lặp từ đầu với bộ nhớ dùng chung: chỉ vòng lặp đầu ghi KV cache đầy đủ, các vòng sau đọc lại cache đó cộng một cửa sổ cục bộ ngắn.
        \item \textbf{Đóng góp nổi bật:}
        \begin{itemize}
            \item \textbf{Chia sẻ bộ nhớ không làm giảm chất lượng mà còn cải thiện nó:} LPT giảm PPL 0,38 đến 1,09 so với Transformer vòng lặp thường, bộ nhớ KV nhỏ hơn 2,0 đến 4,7 lần.
            \item \textbf{Mô hình lai H-LPT:} ở 5 vòng lặp, PPL thấp hơn Transformer chuẩn 1,12 đến 1,82 và dùng ít hơn 76 đến 79\% bộ nhớ ngữ cảnh.
            \item \textbf{Giải thích cơ chế:} đại lộ gradient, phân hóa vai trò giữa các vòng lặp.
        \end{itemize}
    \end{itemize}
\end{frame}

% ------------------------------------------------------------------------------
% Slide 1: Bối cảnh và thách thức
% ------------------------------------------------------------------------------
\begin{frame}{1. LPT: Bối cảnh và cái bẫy bộ nhớ KV cache}
    \textbf{Transformer vòng lặp:} áp dụng cùng một ngăn xếp $L$ tầng $R$ lần cho mỗi token.
    $$h^{(t)} = M_L\big(h_{\text{in}}^{(t)}\big), \qquad h_{\text{in}}^{(t+1)} = N_{\text{rec}}\big(h^{(t)}\big), \qquad t = 1, \dots, R$$
    \begin{itemize}
        \item Độ sâu hiệu dụng $L \cdot R$ nhưng số tham số chỉ bằng $L$ tầng.
    \end{itemize}
    \vspace{0.15cm}
    \textbf{Cái bẫy bộ nhớ:}
    \begin{itemize}
        \item Mỗi vòng lặp ghi một KV cache riêng, nên bộ nhớ KV tăng tuyến tính theo số vòng lặp: $M_{\text{KV}}^{\text{LT}} = L \cdot R \cdot T$.
        \item Ở ngữ cảnh dài, bộ nhớ ngữ cảnh có thể vượt bộ nhớ trọng số, phá hỏng lý do dùng mô hình vòng lặp.
    \end{itemize}
    \vspace{0.15cm}
    \textbf{Giải pháp cũ:} bỏ hoặc dùng lại cache lúc suy luận, không huấn luyện cho việc đó. Chỉ giữ cache của vòng đầu làm chất lượng của Ouro giảm mạnh. Chia sẻ bộ nhớ trước đây chỉ được xem là cách tiết kiệm, chấp nhận mất chất lượng.
\end{frame}

% ------------------------------------------------------------------------------
% Slide 2: Kiến trúc LPT
% ------------------------------------------------------------------------------
\begin{frame}{2. LPT: Kiến trúc Transformer vòng lặp dự đoán}
    \textbf{Tách việc xây bộ nhớ khỏi việc dự đoán theo trục vòng lặp:}
    \begin{itemize}
        \item \textbf{Vòng lặp trạng thái, $t=1$:} chạy như bình thường, ghi KV cache đầy đủ $C^{(1)}$ ở mọi tầng.
        \item \textbf{Vòng lặp dự đoán, $t \ge 2$:} đọc $C^{(1)}$ cộng cửa sổ cục bộ $\omega = 64$ vị trí gần nhất của chính nó, trong cùng một phép softmax:
    \end{itemize}
    $$\mathcal{A}_{\text{LPT}}^{(t)}(i) = C^{(1)}_{\le i} \;\cup\; C^{(t)}_{[i-\omega,\, i]}, \qquad t = 2, \dots, R$$
    \begin{itemize}
        \item \textbf{Bộ nhớ KV:} $M_{\text{KV}}^{\text{LPT}} = L\,[\,T + (R-1)\omega\,]$. Với $T = 4096$, $\omega = 64$, $R = 5$: nhỏ hơn mô hình vòng lặp thường 4,7 lần.
        \item \textbf{Mô hình lai H-LPT:} dùng nền mô hình lai Gated DeltaNet với một tầng chú ý đầy đủ trên năm tầng. Các tầng chú ý tuyến tính giữ trạng thái lặp riêng nhỏ cố định. Các tầng chú ý đầy đủ áp dụng quy tắc trên.
        \item \textbf{Huấn luyện:} mặt nạ khối của FlexAttention trên khóa giá trị dùng chung nối với khóa giá trị cục bộ.
    \end{itemize}
\end{frame}

% ------------------------------------------------------------------------------
% Slide 3: Thiết lập thực nghiệm
% ------------------------------------------------------------------------------
\begin{frame}{3. LPT: Thiết lập thực nghiệm}
    \begin{itemize}
        \item \textbf{Dữ liệu và cỡ mô hình:} FineWeb-Edu, ngữ cảnh 4096. Ba cỡ: S 152M, M 454M, L 1,01B tham số, huấn luyện 20B, 20B và 50B token.
        \item \textbf{Số vòng lặp:} 2, 3, 5. Phụ lục thêm 9.
        \item \textbf{Mô hình đối chứng:} Transformer chuẩn, Transformer sâu mà hẹp khớp độ sâu hiệu dụng và số tham số, Transformer vòng lặp thường, Transformer vòng lặp song song PLT, LOOPIE, mô hình lai LT2.
        \item \textbf{Hai hàm mục tiêu:} mất mát ở vòng cuối, và mất mát thích ứng của Ouro có phân phối dừng học được. Transformer vòng lặp thường dùng mất mát thích ứng, các mô hình còn lại dùng mất mát ở vòng cuối.
        \item \textbf{Đánh giá:} PPL trên 10M token giữ lại, độ chính xác zero-shot trung bình năm tác vụ (ARC-Easy, HellaSwag, PIQA, SciQ, LAMBADA), bộ nhớ và FLOPs tính giải tích tại ngữ cảnh 4096.
        \item \textbf{Không điều chỉnh siêu tham số:} dùng nguyên cấu hình mặc định của nanoGPT cho mọi mô hình.
    \end{itemize}
\end{frame}

% ------------------------------------------------------------------------------
% Slide 4: Kết quả chính
% ------------------------------------------------------------------------------
\begin{frame}{4. LPT: Kết quả chính ở cỡ 1B, 5 vòng lặp}
    \begin{table}[h]
        \centering
        \small
        \begin{tabular}{lccc}
            \toprule
            \textbf{Mô hình} & \textbf{PPL} $\downarrow$ & \textbf{Độ chính xác} $\uparrow$ & \textbf{Bộ nhớ KV so với chuẩn} $\downarrow$ \\
            \midrule
            Transformer chuẩn & 11,35 & 62,5 & $1{,}00\times$ \\
            Transformer vòng lặp thường & 11,39 & 63,4 & $5{,}00\times$ \\
            PLT & 10,59 & 66,4 & $1{,}06\times$ \\
            LT2 mô hình lai & 10,32 & 67,3 & $1{,}01\times$ \\
            \midrule
            \textbf{LPT} & \textbf{10,32} & \textbf{68,2} & $1{,}06\times$ \\
            \textbf{H-LPT} & \textbf{10,23} & 68,1 & $\mathbf{0{,}24\times}$ \\
            \bottomrule
        \end{tabular}
    \end{table}
    \vspace{-0.2cm}
    \begin{small}
    \begin{itemize}
        \item \textbf{LPT so với Transformer vòng lặp thường, mọi cỡ và độ sâu:} PPL thấp hơn 0,38 đến 1,09, độ chính xác cao hơn 0,4 đến 4,8 điểm, bộ nhớ KV nhỏ hơn 2,0 đến 4,7 lần, FLOPs tăng dưới 1\%.
        \item \textbf{H-LPT so với Transformer chuẩn, 5 vòng lặp:} PPL thấp hơn 1,12 đến 1,82, độ chính xác cao hơn 3,1 đến 5,6 điểm, bộ nhớ ngữ cảnh ít hơn 76 đến 79\%.
    \end{itemize}
    \end{small}
\end{frame}

% ------------------------------------------------------------------------------
% Slide 5: Độ sâu và độ ổn định
% ------------------------------------------------------------------------------
\begin{frame}{5. LPT: Mở rộng theo độ sâu và độ ổn định huấn luyện}
    \begin{itemize}
        \item \textbf{Transformer vòng lặp thường ngừng cải thiện từ 3 vòng lặp trở đi}, ở mọi cỡ. LPT tiếp tục cải thiện đến 5 vòng lặp.
        \item \textbf{Hàm mục tiêu ở vòng cuối:}
        \begin{itemize}
            \item Transformer vòng lặp thường huấn luyện bằng mất mát vòng cuối thì xấu dần khi tăng vòng lặp và sụp xuống mức unigram ở 5 vòng, với cấu hình chuẩn hóa của Ouro (S: PPL 2040).
            \item LPT ổn định, và mất mát vòng cuối tốt hơn mất mát thích ứng 0,16 đến 0,50 PPL: các vòng giữa được tự do tinh chỉnh trạng thái mà không phải dự đoán sớm.
        \end{itemize}
        \item \textbf{Chín vòng lặp:} từ 5 lên 9, LPT cải thiện nhẹ ở S nhưng giảm độ chính xác 1,7 điểm ở L. Mô hình lai H-LPT tiếp tục tốt lên ở mọi cỡ.
    \end{itemize}
    \begin{block}{Ý nghĩa}
        Chia sẻ bộ nhớ làm huấn luyện vòng lặp sâu ổn định hơn, nhưng không loại bỏ hoàn toàn bất ổn khi quá sâu.
    \end{block}
\end{frame}

% ------------------------------------------------------------------------------
% Slide 6: Cơ chế 1
% ------------------------------------------------------------------------------
\begin{frame}{6. LPT: Giải mã cơ chế 1 -- Đại lộ gradient}
    \small
    \textbf{Gradient từ mất mát cuối đến vòng lặp đầu yếu đi theo độ sâu:}
    \begin{itemize}
        \item Ở Transformer vòng lặp thường, chuẩn gradient đến vòng đầu giảm mạnh khi tăng số vòng lặp.
        \item Ở LPT, các vòng sau đọc trực tiếp KV cache của vòng đầu, nên gradient có đường tắt về vòng đầu. Ở 5 vòng lặp, phần lớn gradient đến vòng đầu đi qua đường này.
    \end{itemize}
    \vspace{0.1cm}
    \textbf{Bằng chứng nhân quả: cắt gradient của KV cache dùng chung.}
    \begin{itemize}
        \item Hướng thuận không đổi, chỉ cắt đường gradient. PPL tăng ở mọi cấu hình. Cỡ S, 5 vòng: 14,91 tăng lên 33,21. Cỡ M, 3 và 5 vòng: không huấn luyện được, PPL gần mức unigram.
    \end{itemize}
    \vspace{0.1cm}
    \textbf{Điểm cần thận trọng:}
    \begin{itemize}
        \item Tỉ lệ gradient qua KV dùng chung so với gradient qua trạng thái ẩn đã tăng theo độ sâu ngay từ lúc khởi tạo, nên xu hướng này tự nó chưa chứng minh lợi ích.
        \item Tác giả thừa nhận thí nghiệm chưa phân biệt được đường gradient ngắn hơn với thay đổi độ lớn hoặc hướng gradient.
    \end{itemize}
\end{frame}

% ------------------------------------------------------------------------------
% Slide 7: Cơ chế 2
% ------------------------------------------------------------------------------
\begin{frame}{7. LPT: Giải mã cơ chế 2 -- Phân hóa vai trò và cửa sổ cục bộ}
    \textbf{Các vòng dự đoán chú ý chủ yếu vào bộ nhớ dùng chung:}
    \begin{itemize}
        \item Tỉ lệ chú ý vào bộ nhớ dùng chung: 57 đến 71\% ở các vòng dự đoán đầu, 73 đến 80\% ở vòng cuối, kể cả khi cửa sổ cục bộ phủ toàn ngữ cảnh.
    \end{itemize}
    \vspace{0.1cm}
    \textbf{Khóa dùng chung và khóa cục bộ tách thành các cụm riêng trong PCA:}
    \begin{itemize}
        \item Ở Transformer vòng lặp thường, khóa của các vòng chồng lên nhau nhiều hơn.
        \item Vậy sự giống nhau giữa các vòng không phải điều kiện cần để chia sẻ bộ nhớ hiệu quả: vòng đầu chuyên xây bộ nhớ, các vòng sau chuyên dùng nó.
        \item Đây là bằng chứng định tính, trên một chuỗi kiểm tra.
    \end{itemize}
    \vspace{0.1cm}
    \textbf{Quét cửa sổ $\omega$ từ 0 đến 4096:}
    \begin{itemize}
        \item LPT tốt hơn Transformer vòng lặp thường ở mọi $\omega$. Cửa sổ ảnh hưởng ít, điều quyết định là được đọc bộ nhớ dùng chung. Chọn $\omega = 64$.
        \item Lưu ý: ở $\omega = 0$, vòng dự đoán vẫn chú ý vào mục nhập hiện tại của chính nó.
    \end{itemize}
\end{frame}

% ------------------------------------------------------------------------------
% Slide 8: Phép đo GPU
% ------------------------------------------------------------------------------
\begin{frame}{8. LPT: Đo thực tế trên GPU H100}
    \textbf{Bộ nhớ nhỏ hơn cho phép batch lớn hơn.} Ngữ cảnh 4096, bf16, một mã nguồn nghiên cứu dùng chung cho mọi mô hình.
    \begin{itemize}
        \item \textbf{Sức chứa batch:} LPT chứa 93 đến 98\% batch của Transformer chuẩn, và gấp 2,0; 2,9; 4,7 đến 4,8 lần Transformer vòng lặp thường ở 2, 3, 5 vòng lặp. H-LPT chứa gấp 3,6 đến 5,0 lần Transformer chuẩn.
        \item \textbf{Thông lượng của H-LPT so với Transformer chuẩn, tại batch tối đa:} 2,1 đến 2,2 lần ở 2 vòng; 1,4 đến 1,5 lần ở 3 vòng; 0,8 đến 0,9 lần ở 5 vòng.
    \end{itemize}
    \vspace{0.1cm}
    \textbf{Chi phí không biến mất:}
    \begin{itemize}
        \item Cùng batch, LPT chậm hơn Transformer vòng lặp thường 3 đến 8\%, vì vẫn phải đọc cả cache dùng chung. Bộ nhớ dùng chung tiết kiệm dung lượng, không tiết kiệm thời gian đọc.
        \item Batch bằng 1: một bước giải mã của Transformer vòng lặp thường chậm hơn chuẩn 2,3; 3,4; 5,7 lần ở 2, 3, 5 vòng. Ở 5 vòng LPT tốn 4,3 đến 4,8 lần FLOPs của Transformer chuẩn.
    \end{itemize}
\end{frame}

% ------------------------------------------------------------------------------
% Slide 9: Độ tin cậy và công bằng
% ------------------------------------------------------------------------------
\begin{frame}{9. LPT: Độ tin cậy thống kê và tính công bằng}
    \textbf{Bốn hạt giống ngẫu nhiên ở cỡ S, kiểm định t ghép cặp một phía:}
    \begin{table}[h]
        \centering
        \small
        \begin{tabular}{lccc}
            \toprule
            \textbf{Chênh lệch PPL} & \textbf{2 vòng} & \textbf{3 vòng} & \textbf{5 vòng} \\
            \midrule
            LPT so với Transformer vòng lặp thường & $-0{,}40$ & $-0{,}71$ & $-1{,}05$ \\
            H-LPT so với LT2 & $-0{,}04$ & $-0{,}07$ & $-0{,}13$ \\
            LPT so với PLT & $-0{,}07$ & $-0{,}27$ & $-0{,}52$ \\
            \bottomrule
        \end{tabular}
    \end{table}
    \vspace{-0.2cm}
    \begin{small}
    \begin{itemize}
        \item Độ lệch chuẩn giữa các hạt giống lớn nhất là 0,08 PPL. Lợi thế thật sự lớn chỉ khi từ 3 vòng trở lên. Ở 2 vòng, LPT so với PLT gần hòa.
        \item \textbf{Phép so sánh bị lệch có lợi cho LPT:} cấu hình chuẩn hóa của Ouro làm Transformer vòng lặp thường yếu. Thử thêm 5 cách chuẩn hóa khác, bản tốt nhất giảm PPL của nó 0,13 đến 0,76. Khi đó khoảng cách còn \textbf{0,31 đến 0,37} ở S và \textbf{0,06 đến 0,43} ở M.
        \item Chỉ cỡ S có nhiều hạt giống. Cỡ M và L là một lần chạy.
    \end{itemize}
    \end{small}
\end{frame}

% ------------------------------------------------------------------------------
% Slide 10: Hạn chế, phản biện và kết luận
% ------------------------------------------------------------------------------
\begin{frame}{10. LPT: Hạn chế, phản biện và kết luận}
    \textbf{Hạn chế:}
    \begin{itemize}
        \item Không tối ưu về tính toán: 5 vòng lặp tốn 4,3 đến 4,8 lần FLOPs giải mã. Đối chứng khớp tham số, không khớp FLOPs. Mọi token dùng cùng số vòng lặp.
        \item Tối đa 1B tham số, một kho dữ liệu, ngữ cảnh 4096, chưa thử ngữ cảnh dài. Chênh độ chính xác dưới 1 điểm nên xem là nhiễu.
        \item Chưa tối ưu siêu tham số, chưa công bố mã nguồn.
    \end{itemize}
    \vspace{0.1cm}
    \textbf{Kết luận:}
    \begin{itemize}
        \item Chia sẻ bộ nhớ là một thiên kiến quy nạp tốt cho mô hình vòng lặp, không chỉ là mẹo tiết kiệm bộ nhớ.
        \item H-LPT tạo biên chất lượng và bộ nhớ mới cho mô hình vòng lặp: PPL thấp hơn Transformer chuẩn với ít hơn 76 đến 79\% bộ nhớ ngữ cảnh.
        \item \textbf{Hướng tiếp theo:} số vòng lặp thích ứng theo token, nhiều vòng trạng thái hoặc xen kẽ vòng trạng thái và vòng dự đoán, ngữ cảnh dài.
    \end{itemize}
\end{frame}


\begin{frame}
    \centering
    \Huge \textbf{Cảm ơn bạn đã theo dõi!}\\
    \vspace{0.5cm}
    \Large Q\&A
\end{frame}

\end{document}
